% =============================================================================
% 第1章：引言
% =============================================================================

\section{引言}
\label{sec:introduction}

\subsection{研究动机：金融数据隐私与深度学习的张力}

金融反欺诈是商业智能的核心应用场景之一。然而在实际业务中，银行和金融机构之间的数据共享面临严格的隐私合规约束——《个人信息保护法》要求跨机构共享的客户数据须经脱敏处理，GDPR对个人数据的跨境传输施加了明确限制。PCA匿名化（即将原始特征变换为主成分后发布）因其实现简单、数学性质清晰，成为金融数据对外共享中最常见的匿名化手段。例如，Kaggle上广泛使用的信用卡欺诈检测数据集即采用了这一处理——28个原始特征被变换为28个PCA主成分，仅保留交易金额（Amount）作为唯一可识别的业务变量。

这一行业实践带来一个被忽视但至关重要的问题：\textbf{当深度学习模型应用于此类经PCA匿名化处理的金融数据时，其核心机制是否仍能正常工作？}

自注意力机制\cite{vaswani2017attention}是Transformer架构\cite{dosovitskiy2020image}的核心计算原语。在正常条件下，它通过Query-Key内积衡量特征间的语义关联，据此为不同特征分配差异化的注意力权重。然而这一机制的有效运行依赖于一个隐含前提——输入特征具有可区分的语义身份。当特征经过PCA匿名化后，V1--V28等"特征"不再是"交易金额"或"商户类别"这些具有独立业务语义的变量，而是原始特征空间中主成分方向上的坐标值。它们的数学属性高度同构：同样的尺度，同样的分布形态，同样的构造方式。在此条件下，Query-Key内积可能在所有特征对上产生近乎相同的值——注意力机制退化为对所有特征的等权平均，模型的决策过程丧失可区分性和可解释性。

这一问题对金融风控具有直接的商业影响：如果风控模型无法区分"哪个特征更重要"，业务人员就无法理解模型的决策依据，监管合规要求的模型可解释性也无法满足。此外，被等权聚合的特征中若包含与欺诈无关的噪声，将直接损害模型的排序精度。

Gorishniy等人\cite{gorishniy2021revisiting}发现FT-Transformer在部分表格数据集上表现优异，但未系统分析特征匿名化对注意力行为的影响。Grinsztajn等人\cite{grinsztajn2022tree}指出无信息特征对自注意力的干扰可能比树模型更严重。Borisov等人\cite{borisov2024deep}在表格深度学习综述中强调特征结构对注意力有效性的重要性。然而，现有文献均未将"金融数据隐私化→特征匿名化→注意力失效"作为一个完整的因果链来考察。

本文以信用卡欺诈检测为具体场景，聚焦于一个对金融行业有直接操作意义的问题：\textbf{当金融机构使用的数据经过PCA匿名化处理，深度学习模型的自注意力机制是否会系统性失效？若会，失效的根因是什么，有哪些可行的修复路径，各路径的边界在哪里？}

\subsection{研究问题}

\begin{enumerate}
    \item \textbf{诊断问题}：FT-Transformer的自注意力在PCA匿名化金融数据中是否塌缩为均匀分布？根因是特征嵌入的方向同构性，还是Focal Loss对不平衡数据处理的副作用？
    \item \textbf{修复边界问题}：两条修复路径——归一化端（Sparsemax）和数据端（LOR-MIM）——各自的有效范围是什么？联合使用是否存在协同或拮抗？
\end{enumerate}

\subsection{核心发现}

\begin{enumerate}
    \item \textbf{PCA嵌入同构为塌缩根因}：特征嵌入的SVD分析确认29个嵌入向量几乎全部位于3--4个主方向张成的低维子空间，嵌入层有效秩远低于名义维度。
    \item \textbf{两条修复路径各有边界，无模型全面占优}：最优AUPRC/Cost属于Sparsemax（0.9384, 97），最优F1属于LOR-MIM+Softmax（0.9358），最优注意力/ECE属于LOR-MIM+Sparsemax（$\times$16.5, ECE=0.0066）。
    \item \textbf{数据端与归一化端存在结构性拮抗}：LOR-MIM+Sparsemax联合使用使多样性达$\times$16.5，但AUPRC反低于Sparsemax单独使用（0.9267 vs 0.9384）——LOR-MIM的混合破坏了Sparsemax稀疏选择的logits序关系。
\end{enumerate}

\subsection{论文贡献}

\begin{enumerate}
    \item \textbf{商业场景贡献}：明确将金融数据隐私化与深度学习模型失效之间的因果链纳入系统分析，为金融机构在匿名化数据上的模型选型提供了实证依据。
    \item \textbf{诊断贡献}：提供了自注意力塌缩的三层证据链（注意力权重$\to$嵌入SVD$\to$多头多样性）和可操作的诊断模板。
    \item \textbf{方法贡献}：提出LOR-MIM——仅需305参数的数据端对称破缺方法，并揭示其与Sparsemax之间的结构性拮抗关系。
\end{enumerate}

\subsection{论文组织}

第\ref{sec:data}节介绍数据来源及金融匿名化背景；第\ref{sec:background}节回顾FT-Transformer架构、给出塌缩形式化定义并介绍LOR-MIM方法；第\ref{sec:setup}节描述实验设置；第\ref{sec:results}节呈现五变体对照结果；第\ref{sec:diagnosis}节进行深层机制诊断；第\ref{sec:conclusion}节总结全文并给出金融实践建议。
